% B8 · Cold drink — model derivation, step by step
% Compile: pdflatex B8-model-derivation.tex (twice). Overleaf-ready, standard packages only.
\documentclass[11pt,a4paper]{article}

\usepackage[T1]{fontenc}
\usepackage[utf8]{inputenc}
\usepackage{lmodern}
\usepackage[margin=2.2cm]{geometry}
\usepackage{amsmath,amssymb,mathtools}
\usepackage{booktabs,tabularx,array,longtable}
\usepackage{xcolor}
\usepackage[most]{tcolorbox}
\usepackage{enumitem}
\usepackage{microtype}
\usepackage[hidelinks]{hyperref}

\definecolor{accentA}{HTML}{1F5FA8}   % cold / bath
\definecolor{accentB}{HTML}{C2410C}   % warm / drink
\definecolor{ink}{HTML}{1E293B}

\newtcolorbox{result}[1]{enhanced,breakable,colback=accentA!5,colframe=accentA!70,
  boxrule=0.6pt,arc=2pt,left=6pt,right=6pt,top=4pt,bottom=4pt,
  fonttitle=\bfseries\small,coltitle=white,title={#1}}
\newtcolorbox{assume}{enhanced,breakable,colback=accentB!5,colframe=accentB!60,
  boxrule=0.5pt,arc=2pt,left=6pt,right=6pt,top=4pt,bottom=4pt}

\newcommand{\dd}{\mathrm{d}}
\newcommand{\Hb}{\mathcal{H}}
\newcommand{\Qdb}{\dot Q_{db}}
\newcommand{\Tliq}{T_{\mathrm{liq}}}
\newcommand{\wsat}{w_{\mathrm{sat}}}
\newcommand{\Leff}{L_{\mathrm{eff}}}
\newcommand{\Ceff}{C_{\mathrm{eff}}}
\newcommand{\UAbox}{UA_{\mathrm{box}}}
\newcommand{\Tenv}{T_{\mathrm{env}}}
\newcommand{\Tabs}{\mathcal{T}}
\newcommand{\dC}{\,{}^{\circ}\mathrm{C}}
\newcommand{\step}[1]{\section*{#1}\addcontentsline{toc}{section}{#1}}
\newcommand{\sub}[1]{\par\smallskip\noindent\textbf{#1}\enspace}
\newcommand{\means}{\par\noindent\textbf{What it means.}\enspace}

\setlist{itemsep=2pt,topsep=4pt}
\setlength{\parindent}{0pt}
\setlength{\parskip}{5pt}

\title{\textbf{B8 · Cold drink}\\[2pt]\large Model derivation, step by step}
\author{darkpyxlz · SYPT 2027}
\date{25 October 2026}

\begin{document}
\maketitle
\vspace{-1em}

\begin{center}
\begin{minipage}{0.9\linewidth}\small
The model has two parts. The drink is described by one temperature $T_d$. The bath --- ice, brine and
undissolved salt --- is described by its total enthalpy $\Hb$, the heat it holds counted from a
reference state. One heat flow, $UA(T_d - T_b)$, joins them, and a second, through the styrofoam box,
brings heat in from the room. Nothing stirs the bath or the drink.

This document derives the two equations from the conservation of energy and mass and closes them with
the phase diagram of salt water and the laws of heat transfer, explaining every step. Physical
constants and handbook properties of water, ice, salt, brine and the container materials are put in as
numbers, each with a note on what it is (Table~1). Anything that belongs to this rig --- masses, sizes,
starting temperatures, film coefficients, the box's conductance --- stays a symbol, to be measured.
Step~13 collects the results into the closed set of equations; Step~14 solves the cases that have a
formula.
\end{minipage}
\end{center}

\tableofcontents

% ---------------------------------------------------------------
\step{How to read this document}

\textbf{What ``the model'' is.} The 24 Oct salted-ice-bath model, \texttt{b8\_sim.py} in the
repository: a lumped drink; a bath described by its enthalpy, with a map of the five states it can be in
(liquid, ice on the freezing curve, the eutectic plateau, saturated brine with salt, solid); Bodnar's
freezing curve; and the leak of heat through the box. The $\Delta T^{1/4}$ convection law is an option
for the heat path. The buffered eutectic bath (COMSOL v1) is a special case.

\textbf{No stirrer.} Nothing stirs the drink or the bath. Fluid moves only by free convection: liquid
warmer than its surroundings is lighter and rises, colder liquid sinks. That fixes which convection laws
apply (Step~5) and makes the well-mixed bath (assumption A3) the strongest assumption here.

\textbf{How each step is laid out.} It opens with the question it answers and the assumptions it uses.
Each line of algebra is followed by what was done and why. It ends in a boxed \textbf{Result} and a
short note on what the result means.

\textbf{Which quantities become numbers.} A number is put in only for a constant that would be the same
in any laboratory: a physical constant, or a handbook property of a pure substance or a standard
material. These are listed, and explained, in Table~1. The heat capacities, latent heat, heat of solution and
phase diagram are the values \texttt{b8\_sim.py} uses, so the bath formulas here reproduce its numbers. Everything set by this rig stays a symbol (Table~2), because it
has to be measured.

\textbf{Conventions.}
\begin{itemize}
  \item SI units throughout: J, kg, m, s, W. An enthalpy is in J, a heat flow in W.
  \item Temperatures $T$ are in $^\circ$C, because the enthalpy reference and Bodnar's fit are defined
        at $0\dC$. A temperature \emph{difference} is the same in K and $^\circ$C. Absolute temperature,
        in K, appears only in the phase-equilibrium derivation (Step~6), written $\Tabs$.
  \item Salt content $w$ is a mass fraction: kg of NaCl per kg of brine, between 0 and 1. Bodnar's fit
        uses weight per cent, $W = 100w$.
  \item A heat flow is positive in the direction its subscript names: $\Qdb$ flows from the drink to the
        bath, $\dot Q_{\mathrm{leak}}$ from the room to the bath.
  \item Enthalpy reference: liquid water and solid NaCl, both at $0\dC$, have zero enthalpy. Any
        reference would do, because only differences of enthalpy are ever used; this one makes the
        enthalpy of ice at $0\dC$ equal to minus its latent heat.
\end{itemize}

% ---------------------------------------------------------------
\step{Constants put in as numbers}

Each row is a constant that does not depend on this experiment. Where a property changes with
temperature, the value is taken where the model uses it, and the note says how much it changes.

{\small
\begin{longtable}{>{$}l<{$} >{\raggedright\arraybackslash}p{2.7cm} >{\raggedright\arraybackslash}p{9.6cm}}
\caption*{\normalsize\textbf{Table 1.} Constants substituted in this document}\\
\toprule
\text{Symbol} & Value & What it is \\
\midrule
\endfirsthead
\toprule
\text{Symbol} & Value & What it is \\
\midrule
\endhead
\multicolumn{3}{@{}l}{\textit{Physical constants}}\\
g & 9.81\,m\,s$^{-2}$ & Acceleration due to gravity at the Earth's surface. It sets the buoyancy that
  drives free convection. \\
R & 8.314\,J\,mol$^{-1}$K$^{-1}$ & Molar gas constant: Boltzmann's constant times Avogadro's number. It
  converts temperature into energy per mole in the chemical potential. \\
\Tabs_0 & 273.15\,K & Melting point of pure ice at atmospheric pressure: $0\dC$ on the absolute scale. \\
M_w & 0.018015\,kg\,mol$^{-1}$ & Molar mass of water, H$_2$O. \\
M_s & 0.05844\,kg\,mol$^{-1}$ & Molar mass of sodium chloride, NaCl. \\
\nu_i & 2 & Particles each formula unit of NaCl gives in water: Na$^+$ and Cl$^-$. \\
j_{0,1} & 2.405 & First zero of the Bessel function $J_0$; it fixes the slowest cooling pattern of a
  conducting cylinder (Step~3). \\[3pt]
\multicolumn{3}{@{}l}{\textit{Water, ice and salt}}\\
L_f & 333\,550\,J\,kg$^{-1}$ & Latent heat of fusion of ice at $0\dC$: the heat that melts 1\,kg of
  ice at $0\dC$ without warming it. \\
\Delta h_{\mathrm{fus}} & 6009\,J\,mol$^{-1}$ & The same per mole, $L_fM_w$. \\
c_w & 4180\,J\,kg$^{-1}$K$^{-1}$ & Specific heat of liquid water: the heat that warms 1\,kg by 1\,K
  (4180--4220 between $30$ and $0\dC$). \\
c_i & 2000\,J\,kg$^{-1}$K$^{-1}$ & Specific heat of ice. It falls with temperature, from 2097 at $0\dC$
  to 1950 at $-20\dC$ (IAPWS R10-06); 2000 is its value near $-13\dC$, where the bath works. \\
c_h & 864\,J\,kg$^{-1}$K$^{-1}$ & Specific heat of solid NaCl (halite), 50.5\,J\,mol$^{-1}$K$^{-1}$. \\
c_b(w) & $4180 - 4400\,w$ J\,kg$^{-1}$K$^{-1}$ & Specific heat of brine. Dissolved salt lowers it; this
  straight line gives 3300 at $w = 0.20$ and 3160 at $w_E$, 2--5\,\% below the handbook values (3370 and
  3310, Melinder 2010). \\
\Delta h_s & 66\,400\,J\,kg$^{-1}$ & Heat absorbed in dissolving 1\,kg of NaCl in water (3.88\,kJ\,mol$^{-1}$,
  measured at $25\dC$ in dilute solution). Positive: dissolving salt cools the water. \\[3pt]
\multicolumn{3}{@{}l}{\textit{The NaCl--water phase diagram}}\\
b_1,\,b_2,\,b_3 & 1.78, $-0.0442$, $5.572\times10^{-4}$ & Coefficients of Bodnar's (1993) fit to the
  measured freezing curve of salt water, in wt\,\% against degrees below $0\dC$. \\
T_E,\ w_E & $-21.2\dC$, 0.2318 & The eutectic: the coldest brine that can stay liquid, and its salt
  content (the end of Bodnar's fit; handbooks give $-21.1\dC$ and 23.3\,wt\,\%). \\
T_p,\ w_p & $+0.1\dC$, 0.26285 & The peritectic (Swenne 1983): below it the salt that crystallises from
  saturated brine is hydrohalite, NaCl$\cdot$2H$_2$O, not halite. \\
K_f & 1.86\,K\,kg\,mol$^{-1}$ & Cryoscopic constant of water: freezing-point drop per mole of dissolved
  particles per kg of water. Derived in Step~6. \\
\Lambda_E & 241\,450\,J\,kg$^{-1}$ & Heat needed to turn ice and salt at $T_E$ into 1\,kg of eutectic
  brine. Derived in Step~9. \\[3pt]
\multicolumn{3}{@{}l}{\textit{Liquids moving by free convection}}\\
k_d,\,\rho_d & 0.589\,W\,m$^{-1}$K$^{-1}$, 999.1\,kg\,m$^{-3}$ & Thermal conductivity and density of water at
  $15\dC$, the middle of the drink's 28-to-$5\dC$ range. A soft drink is about 90\,\% water. \\
\nu,\ \alpha & $1.139\times10^{-6}$, $1.408\times10^{-7}$\,m$^2$\,s$^{-1}$ & Kinematic viscosity (how
  quickly momentum diffuses) and thermal diffusivity (how quickly heat diffuses) of water at $15\dC$. \\
\beta & $1.51\times10^{-4}$\,K$^{-1}$ & Thermal expansion coefficient of water at $15\dC$: the fractional
  fall in density per kelvin of warming. It is 0 at $4\dC$, $0.88\times10^{-4}$ at 10 and
  $2.07\times10^{-4}$ at $20\dC$. \\
\mathrm{Pr} & 8.1 & Prandtl number of water at $15\dC$, $\nu/\alpha$. \\
k_b,\,\rho_b & 0.54\,W\,m$^{-1}$K$^{-1}$, 1182\,kg\,m$^{-3}$ & Conductivity and density of eutectic brine
  near $-15\dC$ (Melinder 2010). \\
\nu_b,\ \alpha_b & $4.6\times10^{-6}$, $1.38\times10^{-7}$\,m$^2$\,s$^{-1}$ & Kinematic viscosity and
  thermal diffusivity of the same brine; Prandtl number about 33. \\
C_i,\ C_o & 0.62, 0.64 & Constants of laminar free convection on a vertical wall, Churchill \& Chu (1975),
  $0.670/[1 + (0.492/\mathrm{Pr})^{9/16}]^{4/9}$, at $\mathrm{Pr} = 8.1$ (drink) and 33 (brine). \\[3pt]
\multicolumn{3}{@{}l}{\textit{Container and box materials}}\\
c_c & 900\,J\,kg$^{-1}$K$^{-1}$ & Specific heat of aluminium, the can. \\
k_w & 160 or 0.2\,W\,m$^{-1}$K$^{-1}$ & Thermal conductivity of the container wall: can-body aluminium alloy
  (pure aluminium 237), or PET for a plastic bottle. \\
k_{\mathrm{box}} & 0.035\,W\,m$^{-1}$K$^{-1}$ & Thermal conductivity of expanded polystyrene (styrofoam);
  0.033--0.038 depending on the foam's density. \\
\bottomrule
\end{longtable}}

% ---------------------------------------------------------------
\step{Symbols kept as symbols}

Source column: \textbf{M} measured on this rig, \textbf{C} chosen by the experimenter, \textbf{D}
derived here, \textbf{S} a state variable that changes during a run.

{\small
\begin{longtable}{>{$}l<{$} >{\raggedright\arraybackslash}p{8.4cm} l l}
\caption*{\normalsize\textbf{Table 2.} Symbols that stay symbols}\\
\toprule
\text{Symbol} & Meaning & Unit & Src \\
\midrule
\endfirsthead
\toprule
\text{Symbol} & Meaning & Unit & Src \\
\midrule
\endhead
T_d & drink temperature (its well-mixed core) & $^\circ$C & S \\
T_b & bath temperature & $^\circ$C & D (from $\Hb$) \\
\Tenv & room air outside the box & $^\circ$C & M \\
\theta & drink--bath difference, $T_d - T_b$ & K & D \\
m_d,\ c_d & drink mass; drink specific heat (about 4000 for a sugary drink, 4180 for water or a diet drink) & kg, J\,kg$^{-1}$K$^{-1}$ & M \\
m_c & container mass & kg & M \\
C_d & heat capacity of drink + container & J\,K$^{-1}$ & D \\
r_i,\ r_o,\ H & container inner radius, outer radius, height of the wetted wall & m & M \\
t_w & wall thickness & m & M \\
A_i,\ A_o & inner and outer wetted areas & m$^2$ & M \\
h_i,\ h_o & drink-side and bath-side film coefficients & W\,m$^{-2}$K$^{-1}$ & D/M \\
R_i,R_w,R_o,R_c & inner film, wall, outer film and contact (label) resistances & K\,W$^{-1}$ & D \\
U,\ UA & overall heat-transfer coefficient; conductance & W\,m$^{-2}$K$^{-1}$, W\,K$^{-1}$ & M (fit) \\
\mathrm{Bi,Nu,Ra} & Biot, Nusselt and Rayleigh numbers & --- & D \\
\Hb & total enthalpy of the bath & J & S \\
C_t & heat capacity of the tank glass, if counted with the bath & J\,K$^{-1}$ & M \\
m_W,\ m_S & all the water (ice + liquid), all the NaCl & kg & M \\
m_i,\ m_l,\ m_h & ice, brine, undissolved salt & kg & S \\
m_m & mass of ice melted so far & kg & S \\
w,\ \bar w & salt fraction of the brine; overall salt fraction $m_S/(m_W+m_S)$ & --- & S, D \\
\Tliq(w) & freezing temperature of brine of salt fraction $w$ & $^\circ$C & D \\
\wsat(T) & salt fraction of saturated brine at $T$ & --- & D \\
\vartheta & degrees below zero on the freezing curve, $-\Tliq$ & K & D \\
\Leff,\ \Ceff & effective latent heat; effective heat capacity of the bath & J\,kg$^{-1}$, J\,K$^{-1}$ & D \\
C_B & heat capacity of the ice and brine together & J\,K$^{-1}$ & D \\
\UAbox & conductance of the box from room to bath & W\,K$^{-1}$ & M \\
A_{\mathrm{box}},\ \delta_{\mathrm{box}} & foam area and thickness & m$^2$, m & M \\
\Delta\rho & density difference driving the brine film & kg\,m$^{-3}$ & D \\
a_w,\ \varphi,\ b & water activity, osmotic coefficient, molality & ---, ---, mol\,kg$^{-1}$ & D \\
\tau,\ \tau_0,\ t^*,\ \tau_c & time constants & s & D \\
n,\ a & exponent and prefactor in $\Qdb = a\theta^n$ & ---, W\,K$^{-n}$ & D/M \\
T^*,\ T_{fd} & target drink temperature; freezing point of the drink & $^\circ$C & C, M \\
m_{\mathrm{buf}} & salt to keep undissolved as a buffer & kg & C \\
\bottomrule
\end{longtable}}

% ---------------------------------------------------------------
\step{Step 0 --- System, control volumes and assumptions}

\textbf{The question.} What are the parts of the rig, which of them store heat, and along which paths
does heat move between them?

A \emph{control volume} is a region with a boundary drawn round it, kept account of: what crosses the
boundary changes what is inside. The rig splits into two volumes that store heat --- the drink and the
bath --- and two that only resist its flow --- the container wall and the box. Everything outside the
box is the room, so large that the heat it gives up does not change its temperature $\Tenv$.

{\small
\begin{tabularx}{\linewidth}{l X l X}
\toprule
Control volume & Contents & State & Exchanges heat with \\
\midrule
Drink (D) & liquid + container shell & $T_d$ & the bath, through the wall \\
Bath (B) & ice, brine, undissolved NaCl & $\Hb$ & the drink; the room, through the box \\
Container wall & aluminium or PET shell, label & --- & resistance only \\
Box & EPS walls + lid & --- & resistance only \\
\bottomrule
\end{tabularx}}

\medskip
\textbf{Thermal circuit.} Heat flows like current in an electrical circuit: temperatures play the part of
voltages, heat flows of currents, thermal resistances (K\,W$^{-1}$) of electrical ones, and heat
capacities of capacitors, which store heat as their temperature changes.
\[
\underbrace{T_d}_{C_d}
\;\xrightarrow{\;R_i\;}\; T_{wi}
\;\xrightarrow{\;R_c+R_w\;}\; T_{wo}
\;\xrightarrow{\;R_o\;}\;
\underbrace{T_b}_{\Hb}
\;\xleftarrow{\;1/\UAbox\;}\; \Tenv
\]
Heat leaves the drink through the film of drink on the inside of the wall ($R_i$), crosses the label
and the wall ($R_c + R_w$), and enters the bath through the film of brine on the outside ($R_o$). The
wall's two faces are at $T_{wi}$ and $T_{wo}$. From the other side, heat from the room leaks through the
box into the bath. The drink stores heat in $C_d$; the bath stores it in its enthalpy $\Hb$, much of it
as the latent heat of its ice.

\begin{assume}
\textbf{Assumptions,} each with the reason it is made and where it is tested.
\begin{enumerate}[label=\textbf{A\arabic*}]
\item \textbf{Drink well-mixed.} One temperature $T_d$ describes the liquid core; all the change of
  temperature sits in thin boundary layers at the wall. Free convection keeps the core mixed; Step~3
  tests when that is enough.
\item \textbf{Wall and tank store no heat}, or their heat capacity is added to the system they touch:
  the can's to $C_d$, the tank glass's ($C_t$) to the bath. A can wall 0.1\,mm thick holds about 1\,\%
  as much heat per kelvin as its drink.
\item \textbf{Bath well-mixed and at phase equilibrium at every instant.} One temperature $T_b$ and one
  salt fraction $w$ describe the whole bath, and ice, brine and salt are always in equilibrium. With no
  stirrer this is the strongest assumption: brine next to the can warms first. The report and the
  browser simulator relax it, splitting the bath into rings round the can and melting and dissolving at
  finite rates.
\item \textbf{Quasi-steady resistances.} The thin films and the wall reach their steady temperature
  profile much faster (seconds) than $T_d$ or $T_b$ change (minutes), so each is described by a fixed
  resistance.
\item \textbf{Constant pressure, lid on.} Nothing evaporates or condenses, and no frost forms. Radiation
  inside the box is counted in $\UAbox$.
\item \textbf{Constant properties.} Each specific heat, latent heat and heat of solution is one number;
  the brine's specific heat depends on its salt content through one fixed straight line. The
  exceptions are $h_i$ and $\beta$ in the convection laws (Step~5).
\item \textbf{The stable phase diagram, with salt booked as halite.} The freezing curve, the solubility
  and the eutectic are the measured equilibrium ones. Undissolved salt is counted as anhydrous NaCl
  (halite) in the mass and energy balances. Below $+0.1\dC$ the stable salt solid is really hydrohalite,
  NaCl$\cdot$2H$_2$O, which locks up $2M_w/M_s = 0.62$\,kg of water per kg of salt; that bound water and
  its heat are left out. Nothing supercools and no metastable phase forms.
\item \textbf{Drink fully immersed} to a marked waterline, so its wetted area is fixed.
\item \textbf{Ideal sensors.} Probe lag and offsets are corrections to the measurements, not terms in
  the model.
\end{enumerate}
\end{assume}

% ---------------------------------------------------------------
\step{Step 1 --- First law for the drink}

\textbf{The question.} How fast does the drink's temperature fall for a given heat flow out of it?

\sub{1.1 The first law at constant pressure.} The drink and its can form a closed system: no mass
crosses its boundary. For a closed system the first law of thermodynamics says that its internal energy
$E$ changes by the heat that enters it minus the work it does, $\dd E = \delta Q - p\,\dd V$. At constant
pressure (A5) the work of expansion is folded in by using the enthalpy $H = E + pV$, since
$\dd H = \dd E + p\,\dd V$ when $p$ is fixed. So $\dd H = \delta Q$: at constant pressure, \emph{heat in
equals the rise of enthalpy}. The only heat crossing the boundary is $\Qdb$, flowing out, so
\[ \frac{\dd H_D}{\dd t} = -\Qdb . \]

\sub{1.2 The enthalpy of the drink.} The specific heat $c$ is defined by $\dd h = c\,\dd T$ for 1\,kg.
Each kilogram of drink holds $c_d(T_d - T_{\mathrm{ref}})$ above a reference temperature, and each
kilogram of can $c_c(T_d - T_{\mathrm{ref}})$, since the can is at the drink's temperature (A1, A2).
Adding the two masses, with constant specific heats (A6),
\[
H_D = (m_d c_d + m_c c_c)(T_d - T_{\mathrm{ref}}) \equiv C_d\,(T_d - T_{\mathrm{ref}}),
\qquad
\frac{\dd H_D}{\dd t} = C_d\frac{\dd T_d}{\dd t} .
\]
The reference drops out on differentiating. With aluminium's specific heat, $c_c = 900$\,J\,kg$^{-1}$K$^{-1}$:
\[ C_d = m_d c_d + 900\,m_c \qquad (\text{J\,K}^{-1}). \]
The drink's own specific heat $c_d$ stays a symbol because it depends on the drink: sugar lowers it,
from water's 4180 to about 4000\,J\,kg$^{-1}$K$^{-1}$ for a sugary soft drink.

\begin{result}{Result 1 --- drink energy balance}
\[ C_d\,\frac{\dd T_d}{\dd t} = -\Qdb, \qquad C_d = m_d c_d + 900\,m_c \]
\end{result}
\means $C_d$ is the heat the drink and can give up for each kelvin they cool, so the drink cools at
$\Qdb/C_d$ kelvin per second. A can's 13\,g of aluminium adds 12\,J\,K$^{-1}$ to the drink's 1300 or so:
about 1\,\%. Weigh the drink full and empty to get $m_d$; the outside volume $\pi r_o^2H$ would include
the gas in the headspace.

% ---------------------------------------------------------------
\step{Step 2 --- Heat path drink $\to$ bath as series resistances}

\textbf{The question.} How does the heat flow out of the drink depend on the two temperatures, and on
the films and the wall it crosses?

\sub{2.1 Each layer carries the same heat flow.} By A4 the films and the wall store no heat, so whatever
heat enters one face of a layer leaves the other. The same $\Qdb$ therefore crosses all three, each
driven by the temperature drop across it and opposed by its resistance:
\[ \Qdb = \frac{T_d - T_{wi}}{R_i} = \frac{T_{wi}-T_{wo}}{R_w} = \frac{T_{wo}-T_b}{R_o}. \]

\sub{2.2 Film resistances.} Next to a wall, a liquid moving by convection is held still in a thin film,
which heat crosses only by conduction. Newton's law of cooling, $\dot Q = hA\,\Delta T$, sums that up in
a film coefficient $h$, the heat flow per unit area per kelvin of drop. Comparing with
$\dot Q = \Delta T/R$,
\[ R_i = \frac{1}{h_iA_i}, \qquad R_o = \frac{1}{h_oA_o}. \]
The film coefficients depend on how the liquid moves; Step~5 derives them.

\sub{2.3 Wall resistance of a cylinder.} Fourier's law says heat flows down a temperature gradient in
proportion to it: the flux is $-k_w\,\dd T/\dd r$ through a solid of conductivity $k_w$. Through the side
wall of a cylinder of height $H$, it crosses the area $2\pi rH$ at radius $r$, and the total $\dot Q$ is
the same at every radius (no storage, A4):
\begin{gather*}
\dot Q = -k_w(2\pi rH)\frac{\dd T}{\dd r}
\;\Rightarrow\;
\int_{r_i}^{r_o}\frac{\dd r}{r} = -\frac{2\pi k_wH}{\dot Q}\int_{T_{wi}}^{T_{wo}}\dd T
\\
\Rightarrow\;
\ln\frac{r_o}{r_i} = \frac{2\pi k_wH}{\dot Q}\,(T_{wi}-T_{wo})
\;\Rightarrow\;
R_w = \frac{T_{wi}-T_{wo}}{\dot Q} = \frac{\ln(r_o/r_i)}{2\pi k_w H}.
\end{gather*}
The variables were separated ($r$ on one side, $T$ on the other) and both sides integrated across the
wall. For a thin wall, $t_w = r_o - r_i \ll r_i$, the logarithm is $\ln(1 + t_w/r_i)\approx t_w/r_i$, so
$R_w \approx t_w/(k_wA)$ with $A = 2\pi r_iH$: the wall acts as a flat slab. A label sleeve adds a
contact resistance $R_c = R_c''/A$, with $R_c''$ its resistance per unit area, in series.

\sub{2.4 Add the series.} Rearrange each equality in 2.1 as a temperature drop, $T_d - T_{wi} = \Qdb R_i$
and so on, including the label's, and add them. The wall temperatures cancel in pairs:
\[ T_d - T_b = \Qdb\,(R_i + R_c + R_w + R_o). \]
Resistances in series add, because their temperature drops add while the heat flow is common.

\sub{2.5 Faces in parallel.} The side, the base and (if submerged) the top of the container are separate
paths $j$ side by side, each with its own films and wall. Each carries its own heat flow driven by the
same $T_d - T_b$, so the flows add, and with them the conductances (inverse resistances):
\begin{gather*}
UA = \sum_{j}\left(\frac{1}{h_{i,j}A_{i,j}} + R_{c,j} + R_{w,j} + \frac{1}{h_{o,j}A_{o,j}}\right)^{-1},
\\
\frac{1}{U} = \frac{1}{h_i} + R_c'' + \frac{t_w}{k_w} + \frac{1}{h_o}\ \ (\text{thin wall, } A_i\approx A_o\approx A).
\end{gather*}
The second line, for one face with a thin wall, multiplies every resistance by the common area: $U$ is
the overall coefficient per unit area, and $UA$ the conductance in W\,K$^{-1}$.

\sub{2.6 What a probe on the wall measures.} Dividing two of the drops in 2.4 gives the share of the
temperature difference each part takes, as in a potential divider:
\[ \frac{T_d - T_{wo}}{T_{wo}-T_b} = \frac{R_i+R_c+R_w}{R_o}. \]
A probe taped to the outside of the can therefore reads neither the drink nor the bath.

\textbf{Can versus bottle.} With can-body aluminium, $k_w = 160$\,W\,m$^{-1}$K$^{-1}$, a wall 0.1\,mm thick
has $t_w/k_w = 6\times10^{-7}$\,m$^2$K\,W$^{-1}$; a film with $h = 300$ has $1/h = 3\times10^{-3}$, five
thousand times more. So $R_w$ drops out for a can. With PET, $k_w = 0.2$, a wall 0.3\,mm thick has
$t_w/k_w = 1.5\times10^{-3}$, half a film's resistance: a bottle's wall matters. Compare containers by
$U$, not by their cooling time, which also depends on how much drink they hold (Step~4).

\begin{result}{Result 2 --- heat flow from drink to bath}
\[ \Qdb = UA\,(T_d - T_b), \qquad \frac{1}{UA} = R_i + R_c + R_w + R_o \]
\end{result}
\means The whole path from drink to bath is one conductance $UA$, and the heat flow is proportional
to the temperature difference --- as long as the film coefficients do not themselves depend on it,
which Step~5 revisits.

% ---------------------------------------------------------------
\step{Step 3 --- Biot number and when lumping is allowed}

\textbf{The question.} Is it fair to describe the whole drink by one temperature (A1)?

\sub{3.1 The usual test.} Heat leaving a body meets two resistances in series: inside it, across a
length $L_c$ of material of conductivity $k$, about $L_c/k$ per unit area; outside it, $1/U_{\mathrm{ext}}$.
Their ratio is the Biot number,
\begin{gather*}
\mathrm{Bi} = \frac{U_{\mathrm{ext}}L_c}{k_d} = 1.70\,U_{\mathrm{ext}}L_c,
\\
\frac{1}{U_{\mathrm{ext}}} = R_c'' + \frac{t_w}{k_w} + \frac{1}{h_o}, \qquad
L_c = \frac{V}{A}\ \text{or}\ r_i ,
\end{gather*}
with water's conductivity $k_d = 0.589$\,W\,m$^{-1}$K$^{-1}$ ($U_{\mathrm{ext}}$ in W\,m$^{-2}$K$^{-1}$, $L_c$
in m). If $\mathrm{Bi} < 0.1$, the inside resistance is small, the temperature inside hardly varies, and a
body that conducts can be lumped. A 330\,ml can has $V/A \approx 1.3$\,cm; with $U_{\mathrm{ext}}$ of a few
hundred W\,m$^{-2}$K$^{-1}$, $\mathrm{Bi}$ is about 5 to 10. A still drink fails the test. The tiny Biot
number often quoted for a can uses the aluminium's $k_w$ and $t_w$: it belongs to the wall, not the drink.

\sub{3.2 How slowly a still drink would cool.} If the drink only conducted, the heat equation
$\partial T/\partial t = \alpha\nabla^2T$ would hold inside it. In a long cylinder whose wall is held at a
fixed temperature, separating variables gives patterns $J_0(\lambda r)\,e^{-\alpha\lambda^2t}$, where
$J_0$ is a Bessel function; the wall condition $J_0(\lambda r_i) = 0$ allows $\lambda = j_{0,1}/r_i$ for the
slowest pattern, $j_{0,1} = 2.405$. That pattern decays with time constant
\[
\tau_{\mathrm{cond}} = \frac{r_i^2}{\alpha\, j_{0,1}^2} = \frac{r_i^2}{1.408\times10^{-7}\times5.783}
= 1.23\times10^{6}\,r_i^2\ \ \text{s}\qquad(r_i\ \text{in m}),
\]
22 minutes for a can of radius 3.3\,cm: conduction alone would take more than an hour to cool it.

\sub{3.3 The rescue: free convection.} The drink is not still. Liquid cooled at the wall is denser and
sinks in a thin layer down the wall, while the core is stirred by the circulation this sets up. Step~5
shows that the layer is $\delta \sim H\,\mathrm{Ra}_H^{-1/4}$ thick, where $\mathrm{Ra}_H$ is the Rayleigh
number; with water's properties,
\[
\delta \sim \frac{H^{1/4}}{310\,\Delta T_i^{1/4}} = 3.2\times10^{-3}\,H^{1/4}\,\Delta T_i^{-1/4}\ \ \text{m},
\]
with $H$ in m and $\Delta T_i$ in K: about 1\,mm for a can 11.7\,cm tall and a 10\,K drop: far thinner than the radius. All the temperature
change sits in that layer, whose resistance is $R_i$, already in the circuit. The measured speed-up over
conduction is $E = \tau_{\mathrm{cond}}/\tau_{\mathrm{measured}}$.

\begin{result}{Result 3 --- lumping criterion}
\[ \text{A1 holds if }\ \delta \sim H\,\mathrm{Ra}_H^{-1/4} \ll r_i\ \text{ and }\ |T_{\mathrm{top}} - T_{\mathrm{bottom}}| \ll T_d - T_b \quad(\text{not } \mathrm{Bi}<0.1) \]
\end{result}
\means The drink is a mixed core inside a thin cold layer, not a conducting solid. The test is whether
the layer is thin and the core uniform, which a probe at the top and one at the bottom of the drink can
check: cold liquid collects at the bottom, so the core does stratify slowly.

% ---------------------------------------------------------------
\step{Step 4 --- Fixed-bath solution}

\textbf{The question.} If the bath stays at one temperature, how does the drink cool?

\sub{4.1 Combine Results 1 and 2.} Put $\Qdb = UA(T_d - T_b)$ into $C_d\,\dd T_d/\dd t = -\Qdb$. With $T_b$
constant, $\dd\theta/\dd t = \dd T_d/\dd t$ for the difference $\theta = T_d - T_b$, so
\[ \frac{\dd\theta}{\dd t} = -\frac{\theta}{\tau},\qquad \tau \equiv \frac{C_d}{UA}. \]
The difference falls at a rate proportional to itself; $\tau$ (seconds) is the drink's heat capacity over
its conductance.

\sub{4.2 Separate and integrate.} Put everything in $\theta$ on the left and in $t$ on the right, and
integrate from the start, where $\theta = \theta_0$:
\[
\int_{\theta_0}^{\theta}\frac{\dd\theta'}{\theta'} = -\int_0^t\frac{\dd t'}{\tau}
\;\Rightarrow\;
\ln\frac{\theta}{\theta_0} = -\frac{t}{\tau}.
\]

\sub{4.3 Solve for the temperature.} Exponentiate both sides and add back $T_b$: this gives Result~4. In
each time $\tau$ the remaining difference shrinks by the factor $e = 2.718$.

\sub{4.4 Derived quantities.} Setting $T_d = T^*$ and solving for $t$ gives the time to reach a target;
differentiating at $t = 0$ gives the initial cooling rate:
\[
t_{T^*} = \tau\ln\frac{T_{d0}-T_b}{T^*-T_b}, \qquad
\left.\frac{\dd T_d}{\dd t}\right|_0 = -\frac{T_{d0}-T_b}{\tau}.
\]
Two baths at $T_{b1}$ and $T_{b2}$ can be compared in two ways:
\[
\frac{t_{T^*}(T_{b1})}{t_{T^*}(T_{b2})} = \frac{\ln[(T_{d0}-T_{b1})/(T^*-T_{b1})]}{\ln[(T_{d0}-T_{b2})/(T^*-T_{b2})]},
\qquad
\frac{\dot T_d(0)|_{T_{b1}}}{\dot T_d(0)|_{T_{b2}}} = \frac{T_{d0}-T_{b1}}{T_{d0}-T_{b2}} .
\]
The two ratios differ: ``how much faster does the colder bath cool the drink'' has two honest answers,
one for the time to a target and one for the initial rate.

\sub{4.5 Geometry inside $\tau$.} The drink's heat capacity is about $\rho_dc_dV$ for a volume $V$, so
$\tau = (\rho_dc_d/U)\,(V/A)$: at the same $U$, a container with more drink per unit of wetted area cools
more slowly.

\begin{result}{Result 4 --- fixed-bath cooling curve}
\[ T_d(t) = T_b + (T_{d0}-T_b)\,e^{-t/\tau}, \qquad \tau = \frac{C_d}{UA} \]
\end{result}
\means A plot of $\ln(T_d - T_b)$ against time is a straight line of slope $-1/\tau$; from its slope and
the weighed $C_d$, $UA = C_d/\tau$.

% ---------------------------------------------------------------
\step{Step 5 --- Convection closures and the power-law cooling curve}

\textbf{The question.} What sets the film coefficients when nothing is stirred, and how does that change
the cooling curve?

\sub{5.1 Drink side: laminar free convection.} In the layer next to the wall the drink is colder by about
$\Delta T_i = T_d - T_{wi}$, so denser by $\rho\beta\,\Delta T_i$: a buoyancy force $\rho g\beta\,\Delta T_i$
per unit volume drives it down the wall. In a layer of thickness $\delta$ moving at speed $u$, viscosity
resists with a force $\rho\nu u/\delta^2$; balancing the two gives $u \sim g\beta\Delta T_i\delta^2/\nu$.
Heat carried along the wall over its height $H$, $u\Delta T_i/H$, must balance heat conducted across the
layer, $\alpha\Delta T_i/\delta^2$. Eliminating $u$ gives $\delta \sim H\,\mathrm{Ra}_H^{-1/4}$, and since the
film coefficient is about $k_d/\delta$,
\[
\mathrm{Ra}_H = \frac{g\beta\,\Delta T_i H^3}{\nu\alpha}, \qquad
\mathrm{Nu}_H = \frac{h_iH}{k_d} = C_i\,\mathrm{Ra}_H^{1/4} .
\]
The Rayleigh number compares buoyancy with the two diffusions, of momentum and of heat, that resist it;
the Nusselt number compares the film coefficient with conduction over the height. The scaling fixes the
power $\tfrac14$; the constant comes from experiments on vertical walls, $C_i = 0.62$ for water
(Table~1). With water at $15\dC$,
\[
\mathrm{Ra}_H = \frac{9.81\times1.51\times10^{-4}}{1.139\times10^{-6}\times1.408\times10^{-7}}\,\Delta T_iH^3
= 9.22\times10^{9}\,\Delta T_iH^3 .
\]
Solving the Nusselt relation for $h_i$ and putting in the numbers:
\begin{gather*}
h_i = C_ik_d\left(\frac{g\beta}{\nu\alpha H}\right)^{1/4}\Delta T_i^{1/4}
= 0.62\times0.589\times310\;H^{-1/4}\,\Delta T_i^{1/4}
\equiv K_i\,\Delta T_i^{1/4},
\\
K_i = 112\,H^{-1/4}
\end{gather*}
in W\,m$^{-2}$K$^{-5/4}$ with $H$ in m. For an upright can 11.66\,cm tall, $K_i = 192$, and at a 25\,K drop
$h_i = 430$\,W\,m$^{-2}$K$^{-1}$, close to the 400 measured on a bottle. Three remarks. The law holds while
the layer is laminar, $\mathrm{Ra}_H < 10^9$, that is $\Delta T_iH^3 < 0.11$\,m$^3$K: always, here. Sugar
makes a drink about a third more viscous than water, which lowers $K_i$ by under 10\,\%, since it enters
as a quarter power. And for a can lying on its side, as in the rig, the height of the wall is its
diameter, and the constant changes a little (Churchill \& Chu give a separate law for horizontal
cylinders).

\sub{5.2 When the drink's film dominates.} If $R_i$ is much larger than $R_c+R_w+R_o$, almost all the
temperature difference falls across the drink's film, $\Delta T_i \approx \theta$. Then
$\Qdb = h_iA_i\theta = K_iA_i\,\theta^{1/4}\theta$:
\[ \Qdb = K_iA_i\,\theta^{5/4} \equiv a\,\theta^{5/4}. \]
The heat flow grows faster than the difference, because a larger difference also thins the film.

\sub{5.3 Integrate with $T_b$ fixed.} Result~1 becomes $C_d\,\dd\theta/\dd t = -a\theta^{5/4}$. Separate and
integrate as in Step~4, using $\int\theta^{-5/4}\dd\theta = -4\theta^{-1/4}$:
\[
\int_{\theta_0}^{\theta}\theta'^{-5/4}\dd\theta' = -\frac{a}{C_d}\int_0^t\dd t'
\;\Rightarrow\;
\theta^{-1/4} = \theta_0^{-1/4} + \frac{at}{4C_d}
\;\Rightarrow\;
\theta(t) = \theta_0\left(1+\frac{t}{t^*}\right)^{-4}.
\]
The last step factored out $\theta_0^{-1/4}$ and raised both sides to the power $-4$, which defines
\[
t^* = \frac{4C_d}{a\,\theta_0^{1/4}} = 4\tau_0, \qquad \tau_0 = \frac{C_d}{h_i(\theta_0)A_i}.
\]
$\tau_0$ is the time constant the drink would have if its film kept its starting coefficient. The curve
starts like an exponential with that time constant and then slows, as the film thickens.

\sub{5.4 Any exponent.} For $\dot Q = a\theta^n$ with $n\neq1$, the same steps with
$\int\theta^{-n}\dd\theta = \theta^{1-n}/(1-n)$ give
\[
\theta^{1-n} = \theta_0^{1-n} + (n-1)\frac{a}{C_d}t
\;\Rightarrow\;
\theta(t) = \theta_0\big[1+(n-1)t/\tau_0\big]^{-1/(n-1)}, \qquad \tau_0 = \frac{C_d}{a\theta_0^{\,n-1}} .
\]
As $n\to1$, $[1 + (n-1)x]^{-1/(n-1)} \to e^{-x}$ and Result~4 comes back.

\sub{5.5 Testing the exponent.} Taking logarithms of $C_d\,\dd T_d/\dd t = -a(T_d-T_b)^n$,
\[ \ln\!\left(-\frac{\dd T_d}{\dd t}\right) = \ln\frac{a}{C_d} + n\ln(T_d - T_b), \]
so the slope of the log of the cooling rate against the log of the difference is $n$: 1 for a constant
$U$, $\tfrac54$ for a convection-limited drink.

\sub{5.6 Mixed resistances.} When the other resistances are not small, the drink's film sees only part of
$\theta$: what is left after the drop $\Qdb(R_c+R_w+R_o)$ across the rest. That makes $\Qdb$ implicit,
\[ \Qdb = K_iA_i\big[\theta - \Qdb\,(R_c+R_w+R_o)\big]^{5/4}, \qquad 1 \le n_{\mathrm{eff}} \le \tfrac54 , \]
and the effective exponent lies between the two limits.

\sub{5.7 Bath side, unstirred.} The same reasoning applies to the brine on the outside of the can, except
that its density difference $\Delta\rho$ comes from two causes: the brine at the wall is warmer, and the
ice it melts dilutes it. Both make it lighter, so it rises along the wall. Writing the Rayleigh number with
$\Delta\rho$ directly,
\begin{gather*}
\mathrm{Nu}_o = \frac{h_oH}{k_b} = C_o\left(\frac{g\,\Delta\rho\,H^3}{\rho_b\nu_b\alpha_b}\right)^{1/4}
\\
\Rightarrow\;
h_o = C_ok_b\left(\frac{g}{\rho_b\nu_b\alpha_b}\right)^{1/4}\Delta\rho^{1/4}H^{-1/4}
= 118\,\Delta\rho^{1/4}H^{-1/4}
\end{gather*}
in W\,m$^{-2}$K$^{-1}$, with $\Delta\rho$ in kg\,m$^{-3}$ and $H$ in m, using eutectic brine's properties
($C_o = 0.64$, $k_b = 0.54$, $\rho_b = 1182$, $\nu_b = 4.6\times10^{-6}$, $\alpha_b = 1.38\times10^{-7}$).
$\Delta\rho$ stays a symbol: it depends on how warm the wall is and how much meltwater it makes. On a
can 11.7\,cm tall, a few kg\,m$^{-3}$ gives $h_o$ of 250--300\,W\,m$^{-2}$K$^{-1}$.

\sub{5.8 Caveat: $\beta$ changes with temperature.} Water is densest at $4\dC$, where $\beta = 0$; below
that it expands on cooling. Dissolved sugar and CO$_2$ move the maximum a little lower. As the drink
approaches it the buoyancy in its film vanishes, and $\beta\,\Delta T_i$ must be replaced by the density
difference itself:
\[ \mathrm{Ra}_H = \frac{g\,[\rho_d(T_{wi})-\rho_d(T_d)]\,H^3}{\rho_d\,\nu\,\alpha}. \]
Near the maximum the film coefficient falls sharply, so fit $n$ only over the early part of a run.

\begin{result}{Result 5 --- convection closure and fixed-bath power law}
\begin{gather*}
h_i = K_i(T_d - T_{wi})^{1/4},\qquad K_i = 112\,H^{-1/4}\\
\Qdb = a\theta^n,\ 1\le n\le\tfrac54, \qquad \theta(t) = \theta_0\left(1+\frac{t}{t^*}\right)^{-4}\ (n=\tfrac54)
\end{gather*}
\end{result}
\means With nothing stirred, the films set the rate, and each grows as the quarter power of its own
temperature drop. A constant $U$ (Result~2) is the simplest fit to a run; the power law is the
physically expected shape, and the exponent test of 5.5 tells them apart.

% ---------------------------------------------------------------
\step{Step 6 --- Phase equilibrium of the bath}

\textbf{The question.} At what temperature can ice and brine of a given salt content coexist, how much
salt can brine hold, and what is the coldest the bath can get?

\sub{6.1 Equilibrium condition.} Water molecules can pass between the ice and the brine. They stop
moving on balance when water has the same \emph{chemical potential} --- Gibbs energy per mole --- in both.
Ice is pure, so its chemical potential is that of pure ice, $\mu_{\mathrm{ice}}$. In brine the dissolved
ions dilute the water and lower its chemical potential below that of pure liquid water, $\mu_w^*$, by
$R\Tabs\ln a_w$, where $a_w$ is the water's activity ($a_w = 1$ for pure water, less in brine) and
$R = 8.314$\,J\,mol$^{-1}$K$^{-1}$. Equilibrium at absolute temperature $\Tabs$ needs
\[
\mu_{\mathrm{ice}}(\Tabs) = \mu_w^*(\Tabs) + R\Tabs\ln a_w
\;\Rightarrow\;
R\Tabs\ln a_w = -\Delta g_{\mathrm{fus}}(\Tabs),
\]
where $\Delta g_{\mathrm{fus}} = \mu_w^* - \mu_{\mathrm{ice}}$ is the Gibbs energy of melting one mole of
ice.

\sub{6.2 How $\Delta g_{\mathrm{fus}}$ changes with temperature.} The Gibbs--Helmholtz relation,
$\dd(\Delta g/\Tabs)/\dd\Tabs = -\Delta h/\Tabs^2$, links it to the molar heat of melting,
$\Delta h_{\mathrm{fus}} = L_fM_w = 6009$\,J\,mol$^{-1}$. Pure ice and water coexist at
$\Tabs_0 = 273.15$\,K, so $\Delta g_{\mathrm{fus}}(\Tabs_0) = 0$. Integrating from $\Tabs_0$ with
$\Delta h_{\mathrm{fus}}$ constant, and using 6.1,
\begin{gather*}
\frac{\Delta g_{\mathrm{fus}}}{\Tabs} = \Delta h_{\mathrm{fus}}\left(\frac{1}{\Tabs}-\frac{1}{\Tabs_0}\right)
\;\Rightarrow\;
\ln a_w = -\frac{\Delta h_{\mathrm{fus}}}{R}\left(\frac{1}{\Tabs}-\frac{1}{\Tabs_0}\right),
\\
\ln a_w = -722.7\left(\frac{1}{\Tabs}-\frac{1}{273.15}\right)
\end{gather*}
with $\Tabs$ in K. Brine of lower water activity is in equilibrium with ice at a lower temperature. Strictly
$\Delta h_{\mathrm{fus}}$ is not constant: it falls by $(c_w - c_i)M_w = 39$\,J\,mol$^{-1}$ per kelvin below
$0\dC$, about 14\,\% by the eutectic --- one reason the measured curve of 6.4 is used in the model.

\sub{6.3 The dilute limit: Blagden's law.} In dilute brine, $a_w$ is close to the mole fraction of water,
$1 - x$, where $x$ is the mole fraction of dissolved particles, and $\ln(1-x)\approx -x$. Per kg of water
there are $\nu_ib$ moles of particles, $b$ being the molality (moles of NaCl per kg of water), against
$1/M_w$ moles of water, so $x \approx \nu_ibM_w$. Also
$1/\Tabs - 1/\Tabs_0 = (\Tabs_0 - \Tabs)/(\Tabs\Tabs_0) \approx \Delta T_f/\Tabs_0^2$ for a small drop
$\Delta T_f$. Equating,
\begin{gather*}
\Delta T_f = \underbrace{\frac{R\Tabs_0^2M_w}{\Delta h_{\mathrm{fus}}}}_{K_f}\nu_i b,
\\
K_f = \frac{8.314\times273.15^2\times0.018015}{6009} = 1.86\ \text{K\,kg\,mol}^{-1}.
\end{gather*}
The molality of brine of salt fraction $w$ is $b = (w/M_s)/(1-w) = 17.11\,w/(1-w)$\,mol\,kg$^{-1}$: per kg of
brine, $w$\,kg of salt is $w/M_s$ moles, dissolved in $1-w$\,kg of water. With $\nu_i = 2$,
\[ \Delta T_f = 2\times1.86\times17.11\,\frac{w}{1-w} = 63.6\,\frac{w}{1-w}\ \ \text{K}. \]
At 1\,wt\,\% this gives 0.64\,K; the measured drop is 0.59\,K, because the ions are not independent
particles. The osmotic coefficient $\varphi$ ($\approx 0.93$ here) corrects for that, $\ln a_w = -\varphi\nu_ibM_w$
exactly, and then 6.2 gives, for real brine,
\begin{gather*}
\frac{1}{\Tabs_{\mathrm{liq}}} = \frac{1}{\Tabs_0} + \frac{R}{\Delta h_{\mathrm{fus}}}\varphi\,\nu_i\,b\,M_w ,
\\
\frac{1}{\Tabs_{\mathrm{liq}}} = \frac{1}{273.15} + 8.53\times10^{-4}\,\varphi(w)\,\frac{w}{1-w}\ \ \text{K}^{-1}.
\end{gather*}

\sub{6.4 The measured freezing curve (Bodnar 1993).} Rather than tabulate $\varphi$, the model uses a fit to
measured freezing points. With $\vartheta = -\Tliq$ the number of degrees below zero and $W = 100w$,
\[ W = 1.78\,\vartheta - 0.0442\,\vartheta^2 + 5.572\times10^{-4}\,\vartheta^3 . \]
Given $w$, the cubic is solved for $\vartheta$ between 0 and 21.2. Its slope follows from
$\dd W/\dd\vartheta = 1.78 - 0.0884\,\vartheta + 1.672\times10^{-3}\,\vartheta^2$, with $\dd W = 100\,\dd w$ and
$\dd\Tliq = -\dd\vartheta$:
\[ \frac{\dd\Tliq}{\dd w} = -\frac{100}{1.78 - 0.0884\,\vartheta + 1.672\times10^{-3}\,\vartheta^2}\ \ \text{K} . \]
It is negative --- more salt, colder freezing point --- and steepens from 0.56\,K per wt\,\% near $0\dC$ to
1.52\,K per wt\,\% near the eutectic.

\sub{6.5 The phase rule.} Gibbs's phase rule counts how many quantities can still be chosen freely, $F$,
for $C$ components in $P$ phases: $F = C - P + 2$. The components are water and NaCl, $C = 2$, and the
pressure is fixed, which uses up one: $F = 3 - P$.

{\small
\begin{tabularx}{\linewidth}{X c c X}
\toprule
Phases present & $P$ & $F$ & What is fixed \\
\midrule
brine & 1 & 2 & nothing: $T$ and $w$ are both free \\
ice + brine & 2 & 1 & the temperature by the salt: $T = \Tliq(w)$ \\
brine + solid salt & 2 & 1 & the salt by the temperature: $w = \wsat(T)$ \\
ice + brine + solid salt & 3 & 0 & everything: $T=T_E$, $w=w_E$ \\
ice + solid salt & 2 & 1 & $T$ free, below $T_E$ \\
\bottomrule
\end{tabularx}}

With ice, brine and salt all present, nothing is left to choose: the temperature is pinned at the
eutectic, whatever heat comes in, until one of the three phases runs out.

\sub{6.6 Eutectic and solubility.} The eutectic is where the freezing curve meets the solubility curve:
\[
\Tliq(w_E) = T_E,\qquad \wsat(T_E) = w_E,\qquad T_E = -21.2\dC,\ \ w_E = 0.2318 .
\]
Between the eutectic and the peritectic ($+0.1\dC$, $w_p = 0.26285$) the solubility is taken as the straight
line between them, for $-21.2 \le T \le 0.1\dC$:
\[
\wsat(T) \approx w_E + (w_p - w_E)\frac{T - T_E}{T_p - T_E} = 0.2318 + 1.458\times10^{-3}\,(T + 21.2) .
\]
Saturated brine is a little saltier when warmer: 0.0015 per kelvin. At the peritectic it matches the
handbook solubility of halite at $0\dC$, 35.65\,g per 100\,g of water, $w = 0.2628$.

\begin{result}{Result 6 --- phase constraints}
\begin{gather*}
\text{ice + brine: } T_b = \Tliq(w);\qquad \text{brine + salt: } w = 0.2318 + 1.458\times10^{-3}(T_b + 21.2);\\
\text{all three: } T_b = -21.2\dC,\ w = 0.2318
\end{gather*}
\end{result}
\means The phase diagram ties the temperature to the salt content whenever two phases coexist, and pins
both at the eutectic when three do. That is what makes a salted ice bath a thermostat.

% ---------------------------------------------------------------
\step{Step 7 --- Bath mass balances}

\textbf{The question.} How are the amounts of ice, brine and salt related, given that water and salt are
neither created nor destroyed?

Water is in the ice and in the brine; salt is in the brine and in the undissolved crystals. A kilogram of
brine holds $1-w$\,kg of water and $w$\,kg of salt. Counting each substance,
\[
m_W = m_i + (1-w)m_l, \qquad m_S = w\,m_l + m_h, \qquad m_W + m_S = m_i + m_l + m_h ,
\]
the third being the sum of the first two. While no salt is undissolved ($m_h=0$), these give the brine
mass and its salt fraction from the ice alone:
\[
m_l = m_W + m_S - m_i, \qquad w = \frac{m_S}{m_W+m_S-m_i}.
\]
Melting a mass $\dd m_m$ of ice ($\dd m_i = -\dd m_m$) adds as much water to the brine, and the same salt
is spread through more brine. Differentiating,
\[
\frac{\dd m_l}{\dd m_m} = 1,\qquad \frac{\dd w}{\dd m_m} = -\frac{m_S}{m_l^2} = -\frac{w}{m_l}.
\]
Melting always dilutes the brine. Whether the overall salt fraction $\bar w = m_S/(m_W+m_S)$ lies below or
above $w_E$ decides which phase runs out first on the eutectic plateau (Step~9).

\begin{result}{Result 7 --- mass constraints}
\[ m_W = m_i + (1-w)m_l,\quad m_S = w\,m_l + m_h,\quad 0\le m_i\le m_W,\quad 0\le m_h\le m_S \]
\end{result}
\means Two numbers, the water and the salt put in, fix the bath's make-up; the state only decides how
they are shared out between ice, brine and crystals.

% ---------------------------------------------------------------
\step{Step 8 --- Bath enthalpy and the bath energy equation}

\textbf{The question.} How much heat does the bath hold, and how does it change?

\sub{8.1 Specific enthalpies.} Each is counted from the reference (liquid water and solid salt at $0\dC$),
in J\,kg$^{-1}$, with $T$ in $^\circ$C:
\[
h_{\mathrm{ice}} = 2000\,T - 333\,550,\qquad h_{\mathrm{salt}} = 864\,T,\qquad
h_{\mathrm{brine}} = (4180 - 4400\,w)\,T + 66\,400\,w .
\]
Ice at $0\dC$ holds less heat than water at $0\dC$ by the latent heat $L_f$, and cooling it removes
$c_i = 2000$ per kelvin. Salt starts at the reference and has only its sensible heat, $c_h = 864$ per
kelvin. Brine is made from water and salt at the reference: dissolving its $w$\,kg of salt absorbs
$\Delta h_s = 66\,400$ per kg of salt, and it then warms with its own specific heat $c_b(w)$. The heat of
diluting brine further is left out.

The straight line $c_b(w) = (1-w)\,c_w + w\,(c_w - 4400)$ can be read as water at $c_w = 4180$ plus
dissolved salt at $c_w - 4400 = -220$\,J\,kg$^{-1}$K$^{-1}$. The negative number is not a heat capacity of
salt: it is how much dissolved ions lower the heat capacity of the water round them. It has one useful
consequence: at fixed salt, $m_lc_b(w) = 4180\,m_l - 4400\,m_S$, so each kilogram of water added to the
brine adds exactly $c_w$ to its heat capacity.

\sub{8.2 Total enthalpy.} Enthalpy is extensive: the bath's is the sum of mass times specific enthalpy over
its parts, all at the one temperature $T_b$ (A3):
\begin{multline*}
\Hb = m_i(2000\,T_b - 333\,550) + m_l\big[(4180 - 4400\,w)\,T_b + 66\,400\,w\big]\\
+ 864\,m_hT_b\ \ [+\,C_tT_b].
\end{multline*}
The last term is there if the tank glass is counted with the bath (A2).

\sub{8.3 First law for the bath.} The bath is at constant pressure and does no work: as in Step~1, its
enthalpy rises by the heat that enters it, from the drink and from the room. There is no stirrer, so no
other energy enters.
\[ \frac{\dd\Hb}{\dd t} = \Qdb + \dot Q_{\mathrm{leak}}. \]

\sub{8.4 The leak through the box.} The box is a resistance between the room and the bath (A4), with
conductance $\UAbox$ (Step~12):
\[ \dot Q_{\mathrm{leak}} = \UAbox(\Tenv - T_b). \]

\sub{8.5 Energy check.} Adding Result~1, $C_d\,\dd T_d/\dd t = -\Qdb$, the drink's term cancels:
$\dd(C_dT_d + \Hb)/\dd t = \dot Q_{\mathrm{leak}}$. If the box were perfect ($\UAbox = 0$), $C_dT_d + \Hb$ would
stay constant --- what the drink loses, the bath gains. Any computation must pass this check.

\begin{result}{Result 8 --- bath energy equation}
\[ \frac{\dd\Hb}{\dd t} = UA(T_d - T_b) + \UAbox(\Tenv - T_b) \]
\begin{gather*}
\Hb = m_i(2000\,T_b - 333\,550) + m_l\big[(4180 - 4400\,w)\,T_b + 66\,400\,w\big] + 864\,m_hT_b
\end{gather*}
\end{result}
\means The bath's state variable is the heat it holds, not its temperature. Heat in raises $\Hb$; what
that does to $T_b$ depends on whether it warms the bath or melts ice, which Step~9 decides.

% ---------------------------------------------------------------
\step{Step 9 --- Inverting $\Hb$: the regime map}

\textbf{The question.} Given the bath's enthalpy and what was put in, what is its temperature, and how
much ice, brine and salt does it hold?

There are five unknowns, $T_b$, $m_i$, $m_l$, $m_h$ and $w$. The two mass balances (Step~7) and the
enthalpy (Step~8) give three equations; the phase rule (Step~6) supplies the other two from whichever
phases are present (a missing phase supplies its own zero mass). So the state is fixed:
\[ (T_b,\,m_i,\,m_l,\,m_h,\,w) = \mathcal S(\Hb;\,m_W,\,m_S). \]
Which phases are present depends on $\Hb$; the five possibilities are the regimes below, from warm to cold.

\sub{R-I --- liquid only.} No ice, no crystals: $m_i=m_h=0$, $m_l = m_W+m_S$, $w=\bar w$. The enthalpy
is that of the brine alone, which is linear in $T_b$:
\[ T_b = \frac{\Hb - 66\,400\,m_S}{(m_W+m_S)(4180 - 4400\,\bar w)}. \]

\sub{R-II --- ice and brine on the freezing curve.} No crystals ($m_h=0$); this regime exists only if
$\bar w<w_E$. The brine sits on the freezing curve, $T_b = \Tliq(w)$, and Step~7 gives $m_l = m_S/w$ and
$m_i = m_W + m_S - m_S/w$. Putting these into $\Hb$ makes it a function of $w$ alone:
\begin{multline*}
\Hb_{\mathrm{II}}(w) = \Big(m_W+m_S-\frac{m_S}{w}\Big)\big(2000\,\Tliq(w) - 333\,550\big)\\
+ \frac{m_S}{w}(4180 - 4400\,w)\,\Tliq(w) + 66\,400\,m_S .
\end{multline*}
It is solved for $w$ between $\bar w$ and $w_E$. There is exactly one root: adding heat melts ice
($\dd\Hb/\dd m_m = \Leff > 0$, Step~10), and melting dilutes the brine ($\dd w/\dd m_m < 0$), so $\Hb_{\mathrm{II}}$
falls steadily as $w$ rises. Then $T_b = \Tliq(w)$.

\sub{R-III --- the eutectic plateau.} Ice, brine and crystals together: $T_b = -21.2\dC$ and $w = 0.2318$
(Step~6). Only the brine mass is free, and $\Hb$ is linear in it. With no brine at all, everything is ice
and salt at $T_E$:
\[
\Hb_{E0} = m_W(c_iT_E - L_f) + m_Sc_hT_E = -375\,950\,m_W - 18\,317\,m_S .
\]
Each kilogram of eutectic brine is made from $1-w_E$\,kg of ice and $w_E$\,kg of salt at $T_E$; the
enthalpy this takes is the brine's minus theirs, $\Lambda_E = h_{\mathrm{brine}} - (1-w_E)h_{\mathrm{ice}} - w_Eh_{\mathrm{salt}}$:
\[
\Lambda_E = (1-w_E)L_f + w_E\Delta h_s + T_E\big[c_b(w_E) - (1-w_E)c_i - w_Ec_h\big] .
\]
The first term melts 0.768\,kg of ice, 256\,230\,J; the second dissolves 0.232\,kg of salt, 15\,390\,J; the
third, $-21.2\times(3160 - 1536 - 200) = -30\,180$\,J, is negative because the brine has a larger heat
capacity than the ice and salt it came from, and at $T_E$, below the reference, a larger heat capacity
means less enthalpy. In all, $\Lambda_E = 241\,450$\,J\,kg$^{-1}$, and
\[
m_l = \frac{\Hb - \Hb_{E0}}{241\,450},\qquad m_i = m_W - 0.7682\,m_l,\qquad m_h = m_S - 0.2318\,m_l .
\]
Heat moves the bath along the plateau by making or freezing eutectic brine. If $\bar w<w_E$ the salt runs
out first, and the bath leaves into R-II; if $\bar w>w_E$ the ice runs out first, and it leaves into R-IV.

\sub{R-IV --- saturated brine and crystals, no ice.} Only if $\bar w>w_E$. All the water is liquid,
$m_l = m_W/(1-w)$, and the brine is saturated, $w = \wsat(T_b) = 0.2318 + 1.458\times10^{-3}(T_b + 21.2)$; the
rest of the salt is crystals, $m_h = m_S - w\,m_l$. Then $\Hb$ is a function of $T_b$ alone:
\begin{multline*}
\Hb_{\mathrm{IV}}(T_b) = \frac{m_W}{1-w}\big[(4180 - 4400\,w)\,T_b + 66\,400\,w\big]\\
+ \Big(m_S - \frac{w\,m_W}{1-w}\Big)864\,T_b,
\end{multline*}
solved for $T_b$ between $T_E$ and $T_\star$, the temperature at which all the salt would just dissolve,
$\wsat(T_\star)=\bar w$, that is $T_\star = -21.2 + (\bar w - 0.2318)/(1.458\times10^{-3})$.

\sub{R-V --- fully solid.} No brine ($m_l=0$): all the water is ice and all the salt crystals, below $T_E$.
The enthalpy is linear in $T_b$:
\[ T_b = \frac{\Hb + 333\,550\,m_W}{2000\,m_W + 864\,m_S}. \]

\textbf{Regime boundaries.} Each threshold is the enthalpy at which a phase appears or disappears, found by
putting that state into the formula of either neighbouring regime; the state is continuous across every
boundary.

{\small
\begin{tabularx}{\linewidth}{l l l X}
\toprule
Bath type & Regime & Range & Threshold \\
\midrule
$\bar w<w_E$ & R-I & $\Hb\ge\Hb_A$ & $\Hb_A = (m_W+m_S)(4180 - 4400\,\bar w)\Tliq(\bar w) + 66\,400\,m_S$: first ice \\
 & R-II & $\Hb_B\le\Hb<\Hb_A$ & $\Hb_B = \Hb_{E0} + 1\,041\,600\,m_S$: all salt in eutectic brine \\
 & R-III & $\Hb_{E0}\le\Hb<\Hb_B$ & \\
 & R-V & $\Hb<\Hb_{E0}$ & \\
\midrule
$\bar w>w_E$ & R-I & $\Hb\ge\Hb_C$ & $\Hb_C = (m_W+m_S)(4180 - 4400\,\bar w)T_\star + 66\,400\,m_S$: first crystal \\
 & R-IV & $\Hb_D\le\Hb<\Hb_C$ & $\Hb_D = \Hb_{E0} + 314\,300\,m_W$: all water in eutectic brine \\
 & R-III & $\Hb_{E0}\le\Hb<\Hb_D$ & \\
 & R-V & $\Hb<\Hb_{E0}$ & \\
\bottomrule
\end{tabularx}}

The numbers are $\Lambda_E/w_E = 1\,041\,600$ and $\Lambda_E/(1-w_E) = 314\,300$\,J\,kg$^{-1}$: the enthalpy
of the brine that holds all the salt, or all the water, at the eutectic.

\begin{result}{Result 9 --- bath closure}
\[ T_b = T_b(\Hb),\quad m_i = m_i(\Hb),\quad m_h = m_h(\Hb),\quad w = w(\Hb) \quad\text{via R-I to R-V} \]
\end{result}
\means The enthalpy alone, with the two masses, tells everything about the bath. This is why the model
steps $\Hb$ forward in time rather than $T_b$: $\Hb$ always changes smoothly, while $T_b$ stops dead on the
plateau.

% ---------------------------------------------------------------
\step{Step 10 --- Effective latent heat, bath drift and the blank run}

\textbf{The question.} On the freezing curve, how much heat melts a kilogram of ice, and how fast does the
bath warm?

\sub{10.1 Differentiate along the freezing curve.} In R-II all the salt is dissolved, so
$m_lw = m_S$ is fixed, and by 8.1 $m_l(4180 - 4400w) = 4180\,m_l - 4400\,m_S$. The enthalpy is
$\Hb = m_i(c_iT_b - L_f) + (4180\,m_l - 4400\,m_S)T_b + 66\,400\,m_S$, with $m_l = m_W+m_S-m_i$. A small change
of ice and temperature changes it by
\[
\dd\Hb = \big[h_{\mathrm{ice}}(T_b) - c_wT_b\big]\dd m_i
+ \underbrace{\big(m_ic_i + m_lc_b(w)\big)}_{C_B}\dd T_b .
\]
The first term: a kilogram of ice that melts becomes a kilogram of water in the brine, at heat content
$c_wT_b$ in place of $h_{\mathrm{ice}}(T_b)$. The second warms the ice and brine already there; $C_B$ is
their heat capacity.

\sub{10.2 Heat per kilogram melted.} With $\dd m_m = -\dd m_i$ and $h_{\mathrm{ice}} = c_iT_b - L_f$:
\begin{gather*}
\Leff \equiv \frac{\dd\Hb}{\dd m_m} = L_f + (c_w - c_i)T_b + C_B\frac{\dd T_b}{\dd m_m},
\\
\Leff = 333\,550 + 2180\,T_b + C_B\frac{\dd T_b}{\dd m_m} .
\end{gather*}
Three parts: the latent heat at $0\dC$; a correction because the ice melts at $T_b$, not $0\dC$ (below
zero the latent heat is smaller by $c_w - c_i = 2180$\,J\,kg$^{-1}$ per kelvin, Kirchhoff's law, so this
term is negative); and the heat that warms
the whole bath, because melting dilutes the brine and so raises its freezing point. The last part makes
$\Leff$ larger than $L_f$: for the 20 Oct rig-note charge at $-17.8\dC$, $\Leff = 411\,500$\,J\,kg$^{-1}$,
23\,\% more.

\sub{10.3 Warming per kilogram melted.} By the chain rule, with $\dd w/\dd m_m = -w/m_l$ (Step~7) and the
slope of the freezing curve (6.4),
\[
\frac{\dd T_b}{\dd m_m} = \frac{\dd\Tliq}{\dd w}\,\frac{\dd w}{\dd m_m}
= \frac{100\,w}{m_l\,(1.78 - 0.0884\,\vartheta + 1.672\times10^{-3}\,\vartheta^2)} > 0 .
\]
Melting always warms a bath on the freezing curve, and more so the less brine there is.

\sub{10.4 Effective heat capacity in every regime.} The heat that raises the bath's temperature by one
kelvin, $\Ceff = \dd\Hb/\dd T_b$, follows from differentiating each regime's enthalpy:
\[
\Ceff =
\begin{cases}
(m_W+m_S)(4180 - 4400\,\bar w) & \text{R-I}\\[4pt]
C_B + \dfrac{333\,550 + 2180\,T_b}{\dd T_b/\dd m_m} & \text{R-II}\\[10pt]
\infty & \text{R-III}\\[4pt]
m_lc_b(w) + 864\,m_h + \dfrac{1.458\times10^{-3}\,m_W}{(1-w)^2}\,(66\,400 - 1084\,T_b) & \text{R-IV}\\[10pt]
2000\,m_W + 864\,m_S & \text{R-V}
\end{cases}
\]
R-I and R-V are plain heat capacities. In R-II, warming by one kelvin needs $1/(\dd T_b/\dd m_m)$\,kg of
ice melted, each costing $L_f + (c_w - c_i)T_b$, on top of $C_B$: this is usually much larger than $C_B$.
On the plateau no amount of heat changes the temperature. In R-IV, warming dissolves more salt,
$\dd(w\,m_l)/\dd T_b = m_W\wsat'/(1-w)^2$ with $\wsat' = 1.458\times10^{-3}$\,K$^{-1}$, each kilogram
absorbing $\Delta h_s$ and changing the sensible heat by $(-220 - 864)T_b = -1084\,T_b$ (dissolved instead
of solid salt, 8.1).

\sub{10.5 The plateau as a thermostat.} On the plateau the heat coming in, $\dot Q_{\mathrm{in}} = \Qdb +
\dot Q_{\mathrm{leak}}$, does not warm the bath; it makes eutectic brine at $\Lambda_E$ per kilogram, using ice
and salt in the ratio $1-w_E$ to $w_E$:
\[
\frac{\dd m_l}{\dd t} = \frac{\dot Q_{\mathrm{in}}}{241\,450},\qquad
\frac{\dd m_i}{\dd t} = -\frac{\dot Q_{\mathrm{in}}}{314\,300},\qquad
\frac{\dd m_h}{\dd t} = -\frac{\dot Q_{\mathrm{in}}}{1\,041\,600}\ \ (\text{kg\,s}^{-1}).
\]
The undissolved salt $m_h(0)$ is the buffer: it lasts
\[ t_{\mathrm{buffer}} = \frac{m_h(0)\,\Lambda_E}{w_E\,\overline{\dot Q}_{\mathrm{in}}} = \frac{1.042\times10^{6}\,m_h(0)}{\overline{\dot Q}_{\mathrm{in}}}\ \ \text{s}, \]
with $m_h(0)$ in kg and the mean heat flow in W. Each 100\,g of spare salt holds the bath at $-21.2\dC$
against 1\,W for 29 hours.

\sub{10.6 The blank run.} With no drink in the bath and no stirrer, only the leak brings heat in. On the
freezing curve, $\dot Q_{\mathrm{leak}} = \Leff\,\dd m_m/\dd t$, and since $\dd m_m/\dd t = \dot T_b/(\dd T_b/\dd m_m)$,
\[
\dot Q_{\mathrm{leak}} = \Leff\frac{\dd T_b/\dd t}{\dd T_b/\dd m_m} = \Ceff\frac{\dd T_b}{\dd t}
\;\Rightarrow\;
\UAbox = \frac{\Ceff\,\dot T_b}{\Tenv - T_b}.
\]
A logged warming rate $\dot T_b$ therefore measures the box. Using $L_f$ in place of $\Leff$ would miss the
heat that warms the whole bath and underestimate the leak.

\begin{result}{Result 10 --- explicit bath temperature equation}
\[ \Ceff(T_b)\frac{\dd T_b}{\dd t} = UA(T_d - T_b) + \UAbox(\Tenv - T_b) \qquad(\dot T_b = 0 \text{ in R-III}) \]
\end{result}
\means The bath behaves like a body of heat capacity $\Ceff$, but one that changes with the regime: huge on
the freezing curve, infinite on the plateau. That is why a bath with ice in it warms so slowly.

% ---------------------------------------------------------------
\step{Step 11 --- Initial state from premixing}

\textbf{The question.} What state does the bath start in, once the ice, water and salt are mixed?

\sub{11.1 Enthalpy of the ingredients.} Before mixing, each ingredient is at its own temperature: ice at
$T_{i0}$, water at $T_{w0}$, salt at $T_{s0}$. Using 8.1 for each,
\[ \Hb_0 = m_{i0}(2000\,T_{i0} - 333\,550) + 4180\,m_{w0}T_{w0} + 864\,m_ST_{s0}, \]
and all the water is $m_W = m_{i0}+m_{w0}$.

\sub{11.2 Adiabatic mixing.} In the insulated box, mixing exchanges heat only among the ingredients, so
the enthalpy does not change (Step~8 with no heat in). The starting state is the one Step~9 gives for
$\Hb_0$: $(T_{b0}, m_i, m_l, m_h, w)(0) = \mathcal S(\Hb_0; m_W, m_S)$. Enthalpy is a state function, so
this does not depend on the order of mixing: adding the salt at once, in stages, or pre-dissolved reaches
the same equilibrium, only at different rates.

\sub{11.3 Start on the eutectic.} The bath starts on the plateau if $\Hb_0$ falls in its range (Step~9):
\begin{gather*}
\Hb_{E0}\le\Hb_0<\Hb_{E0} + 1\,041\,600\,m_S\quad\text{if } \bar w<w_E,
\\
\Hb_{E0}\le\Hb_0<\Hb_{E0} + 314\,300\,m_W\quad\text{if } \bar w>w_E.
\end{gather*}

\sub{11.4 Keep a salt buffer.} To start with at least $m_{\mathrm{buf}}$ of salt undissolved, the plateau
formula $m_h = m_S - 0.2318\,m_l$ with $m_l = (\Hb_0 - \Hb_{E0})/241\,450$ must give $m_h \ge m_{\mathrm{buf}}$:
\[ \Hb_0 \le \Hb_{E0} + 1\,041\,600\,(m_S - m_{\mathrm{buf}}). \]

\sub{11.5 Minimum ice fraction.} Write the ice as a fraction $f = m_{i0}/m_W$ of the water. Then
$\Hb_0 = m_W[f(2000\,T_{i0} - 333\,550) + (1-f)\,4180\,T_{w0}] + 864\,m_ST_{s0}$, which falls as $f$ rises.
For $\Hb_0$ to stay below any bound $\Hb_{\max}$ (from 11.3 or 11.4),
\[ f \ge f_{\mathrm{crit}} = \frac{4180\,m_WT_{w0} + 864\,m_ST_{s0} - \Hb_{\max}}{m_W\,(4180\,T_{w0} - 2000\,T_{i0} + 333\,550)}. \]

\sub{11.6 Melting against dissolving.} When the mixture cools itself, the heat comes from melting
$\Delta m$ of ice and dissolving $m_{S,\mathrm{dis}}$ of salt. The share due to melting is
\[ \chi_{\mathrm{lat}} = \frac{333\,550\,\Delta m}{333\,550\,\Delta m + 66\,400\,m_{S,\mathrm{dis}}}. \]
Per kilogram, melting ice absorbs five times as much heat as dissolving salt.

\begin{result}{Result 11 --- initial conditions}
\begin{gather*}
T_d(0) = T_{d0},\\
\Hb(0) = m_{i0}(2000\,T_{i0} - 333\,550) + 4180\,m_{w0}T_{w0} + 864\,m_ST_{s0}
\end{gather*}
\end{result}
\means How cold the bath starts, and whether it starts on the plateau with salt to spare, is set by a
single number, the enthalpy of what goes in; a drier, colder charge has less of it.

% ---------------------------------------------------------------
\step{Step 12 --- The box leak}

\textbf{The question.} How fast does heat from the room get into the bath?

Heat from the room crosses three layers in series: the film of room air on the outside of the box, the
foam, and the air between the foam and the tank. As in Step~2 their resistances add:
\[
\frac{1}{\UAbox} = \frac{1}{h_{\mathrm{air}}A_{\mathrm{out}}} + \frac{\delta_{\mathrm{box}}}{0.035\,A_{\mathrm{box}}}
+ \frac{1}{h_{\mathrm{in}}A_{\mathrm{in}}}
\approx \frac{\delta_{\mathrm{box}}}{0.035\,A_{\mathrm{box}}},
\]
with expanded polystyrene's conductivity, 0.035\,W\,m$^{-1}$K$^{-1}$. The foam dominates: 30\,mm of it has
$\delta/k = 0.86$\,m$^2$K\,W$^{-1}$, against about $1/9 = 0.11$ for room air (free convection plus radiation)
and $1/6 = 0.17$ for the air gap inside. Dropping the films overestimates the leak by about a third;
gaps round the lid, where heat bypasses the foam, push it the other way.

Two variants for when an assumption is broken. With the lid off (A5), the bath's open surface also takes
heat from the air and from water condensing on it: $\dot Q_{\mathrm{surf}} = h_sA_s(T_{\mathrm{air}} - T_b)
+ \dot m_vh_{sg}$, with $\dot m_v$ the condensation rate and $h_{sg}$ its latent heat. With the drink not
fully immersed (A8), its exposed part takes heat from the air:
$C_d\dot T_d = -UA(T_d - T_b) + UA_{da}(T_{\mathrm{air}} - T_d)$.

\begin{result}{Result 12 --- the leak}
\[ \dot Q_{\mathrm{leak}} = \UAbox(\Tenv - T_b),\qquad \UAbox\approx\frac{0.035\,A_{\mathrm{box}}}{\delta_{\mathrm{box}}} \]
\end{result}
\means The leak is proportional to how much colder the bath is than the room, so it is largest exactly
when the bath is most useful. It is best measured, with the blank run of 10.6.

% ---------------------------------------------------------------
\step{Step 13 --- The closed model}

\textbf{The question.} What is the complete set of equations to integrate?

Two first-order differential equations for the state $(T_d, \Hb)$, closed by one law for the heat flow and
one map from $\Hb$ to the bath's state.

\sub{13.1 Governing equations} (Results 1 and 8):
\[
C_d\frac{\dd T_d}{\dd t} = -\Qdb, \qquad
\frac{\dd\Hb}{\dd t} = \Qdb + \UAbox(\Tenv - T_b) .
\]

\sub{13.2 Heat-flow closure --- choose one} (Results 2 and 5):
\begin{gather*}
\text{A (constant conductance):}\quad \Qdb = UA(T_d - T_b), \quad \frac{1}{UA} = R_i+R_c+R_w+R_o\\
\text{B (drink-film convection):}\quad \Qdb = 112\,H^{-1/4}A_i\big[(T_d - T_b) - \Qdb(R_c+R_w+R_o)\big]^{5/4}
\end{gather*}

\sub{13.3 Bath closure} (Results 6, 7 and 9), $(T_b, m_i, m_l, m_h, w) = \mathcal S(\Hb; m_W, m_S)$, in J,
kg and $^\circ$C:
{\small
\[
\mathcal S:\
\begin{cases}
\text{R-I:} & m_i=m_h=0,\ w=\bar w,\ T_b = \dfrac{\Hb - 66\,400\,m_S}{(m_W+m_S)(4180 - 4400\,\bar w)}\\[10pt]
\text{R-II:} & m_h=0,\ m_l = m_S/w,\ T_b=\Tliq(w),\ w \text{ from } \Hb = \Hb_{\mathrm{II}}(w)\ \text{(Step 9)}\\[6pt]
\text{R-III:} & T_b=-21.2,\ w=0.2318,\ m_l = \dfrac{\Hb-\Hb_{E0}}{241\,450},\ m_i = m_W-0.7682\,m_l,\ m_h = m_S - 0.2318\,m_l\\[10pt]
\text{R-IV:} & m_i = 0,\ w=\wsat(T_b),\ T_b \text{ from } \Hb = \Hb_{\mathrm{IV}}(T_b)\ \text{(Step 9)}\\[6pt]
\text{R-V:} & m_l=0,\ T_b = \dfrac{\Hb + 333\,550\,m_W}{2000\,m_W + 864\,m_S}
\end{cases}
\]}
with the regime chosen by the thresholds of Step~9, and
\begin{gather*}
W = 1.78\,\vartheta - 0.0442\,\vartheta^2 + 5.572\times10^{-4}\,\vartheta^3\ \ (W=100w,\ \vartheta=-\Tliq),\\
\wsat(T) = 0.2318 + 1.458\times10^{-3}(T + 21.2),\qquad
\Hb_{E0} = -375\,950\,m_W - 18\,317\,m_S,\qquad
\bar w = \frac{m_S}{m_W+m_S}.
\end{gather*}

\sub{13.4 Equivalent explicit form.} For a hand calculation (Desmos, Mathematica), Result~10 replaces the
$\Hb$ equation:
\[
C_d\dot T_d = -\Qdb, \qquad \Ceff(T_b)\,\dot T_b = \Qdb + \UAbox(\Tenv - T_b)\quad(\dot T_b = 0 \text{ in R-III}).
\]
A program should integrate the $\Hb$ form, which crosses the plateau without dividing by $\Ceff=\infty$.

\sub{13.5 Initial conditions} (Result 11):
\begin{gather*}
T_d(0) = T_{d0},\qquad m_W = m_{i0}+m_{w0},\\
\Hb(0) = m_{i0}(2000\,T_{i0}-333\,550) + 4180\,m_{w0}T_{w0} + 864\,m_ST_{s0}.
\end{gather*}

\sub{13.6 When to stop, and when the model stops being valid.}
\begin{itemize}
  \item $T_d$ reaches the target $T^*$: the measured output.
  \item $T_d$ reaches the drink's freezing point $T_{fd}$ (about $-0.9\dC$ for a sugary fizzy drink):
        beyond it the drink needs its own enthalpy model, like the bath's.
  \item In a buffered run, $m_h\to0$: the thermostat is gone and the bath leaves the plateau.
  \item $T_d$ approaches the density maximum near $4\dC$: closure B and assumption A1 weaken (5.8).
\end{itemize}

\sub{13.7 What has to be measured.}

{\small
\begin{tabularx}{\linewidth}{l X X}
\toprule
Group & Parameters & How they are fixed \\
\midrule
Drink & $m_d$, $c_d$, $m_c$ & weighing; the drink's specific heat from its label (sugar) \\
Heat path & $UA$ (A), or $A_i$, $H$, $R_c$, $R_w$, $R_o$ (B) & fitted on a buffered run, where $T_b = -21.2\dC$ is known \\
Bath & $m_W$, $m_S$, $m_{i0}$, $m_{w0}$ & weighing \\
Environment & $\UAbox$, $\Tenv(t)$ & a blank run (10.6); room temperature logged on channel 4 \\
Start & $T_{d0}$, $T_{i0}$, $T_{w0}$, $T_{s0}$ & logged \\
\bottomrule
\end{tabularx}}

Every other constant is in Table~1. Once $UA$ is fixed on a buffered run, an unbuffered run has no free
parameter left: it is a prediction.

\begin{tcolorbox}[enhanced,colback=accentB!6,colframe=accentB!80,boxrule=1pt,arc=2pt,
  fonttitle=\bfseries,coltitle=white,title={Result 13 --- the model, in one block}]
\[
\begin{aligned}
C_d\,\dot T_d &= -\,UA\,(T_d - T_b), \qquad C_d = m_dc_d + 900\,m_c\\
\dot{\Hb} &= UA\,(T_d - T_b) + \UAbox(\Tenv - T_b)\\
T_b &= T_b(\Hb;\,m_W,\,m_S)\ \text{ via } \mathcal S\\
T_d(0) &= T_{d0},\qquad \Hb(0) = \Hb_0
\end{aligned}
\]
\end{tcolorbox}
\means Five numbers describe the rig ($C_d$, $UA$, $m_W$, $m_S$, $\UAbox$), plus the starting
temperatures; the rest is physics that is the same in any laboratory.

% ---------------------------------------------------------------
\step{Step 14 --- Closed-form special cases}

\textbf{The question.} When can the equations be solved with a formula, rather than step by step?

\sub{14.1 Buffered bath.} If the bath stays on the plateau (R-III throughout), $T_b = -21.2\dC$ is constant
and Result~4 applies directly:
\[ T_d(t) = -21.2 + (T_{d0} + 21.2)\,e^{-t/\tau}, \qquad \tau = \frac{C_d}{UA}, \]
valid while $t < t_{\mathrm{buffer}}$ (10.5).

\sub{14.2 Bath on the freezing curve, sealed box.} In R-II over a small range of temperature, $\Ceff$ is
roughly constant; take $\UAbox = 0$. Then drink and bath are two heat capacities joined by one
conductance. Dividing the two equations of 13.4 by their heat capacities and subtracting,
\[
\frac{\dd\theta}{\dd t} = -UA\Big(\frac{1}{C_d}+\frac{1}{\Ceff}\Big)\theta
\;\Rightarrow\;
\theta = \theta_0 e^{-t/\tau_c},\qquad \tau_c = \frac{C_d\Ceff}{UA\,(C_d+\Ceff)} = \frac{\tau}{1+\Gamma},\quad \Gamma = \frac{C_d}{\Ceff}.
\]
The difference closes faster than with a fixed bath, because the bath warms to meet the drink. Energy
conservation (8.5), $C_dT_d + \Ceff T_b = \text{const}$, fixes the common final temperature, and each
temperature approaches it with the same time constant:
\[
T_\infty = \frac{C_dT_{d0} + \Ceff T_{b0}}{C_d+\Ceff},\qquad
T_d(t) = T_\infty + (T_{d0}-T_\infty)e^{-t/\tau_c},\qquad
T_b(t) = T_\infty - (T_\infty - T_{b0})e^{-t/\tau_c}.
\]
The drink reaches $T^*$ only if $T_\infty < T^*$, after $t_{T^*} = \tau_c\ln\big[(T_{d0}-T_\infty)/(T^*-T_\infty)\big]$.

\sub{14.3 With the box leak.} With $\Tenv$ constant, the two equations of 13.4 are linear in
$\mathbf x = (T_d, T_b)^{\mathsf T}$: $\dot{\mathbf x} = \mathsf A\mathbf x + \mathbf f$, with
$\mathbf f = (0,\ \UAbox\Tenv/\Ceff)^{\mathsf T}$ and
\[
\mathsf A = \begin{pmatrix} -\dfrac{UA}{C_d} & \dfrac{UA}{C_d}\\[8pt] \dfrac{UA}{\Ceff} & -\dfrac{UA+\UAbox}{\Ceff}\end{pmatrix}.
\]
The solution is a steady state plus two decaying exponentials, $e^{\lambda_\pm t}$, where $\lambda_\pm$ are the
eigenvalues of $\mathsf A$, the roots of $\lambda^2 - (\operatorname{tr}\mathsf A)\lambda + \det\mathsf A = 0$:
\begin{gather*}
\lambda_\pm = \tfrac12\Big[\operatorname{tr}\mathsf A \pm \sqrt{(\operatorname{tr}\mathsf A)^2 - 4\det\mathsf A}\Big],
\\
\operatorname{tr}\mathsf A = -\Big(\frac{UA}{C_d} + \frac{UA+\UAbox}{\Ceff}\Big),\qquad
\det\mathsf A = \frac{UA\cdot\UAbox}{C_d\,\Ceff}.
\end{gather*}
The trace is negative and the determinant positive, so both eigenvalues are real and negative. When the
leak is small, the fast one is nearly $\operatorname{tr}\mathsf A$: the drink and bath coming together as
in 14.2. The slow one is nearly $\det\mathsf A/\operatorname{tr}\mathsf A$: the room warming both. The steady state is $T_d = T_b = \Tenv$, reached only after the ice is gone --- by which time
$\Ceff$ has changed, so this solution holds only within one regime.

\sub{14.4 Fixed bath, convection-limited drink.} Result~5: $\theta = \theta_0(1+t/t^*)^{-4}$, $t^* = 4\tau_0$.

\means Each special case isolates one physical effect: the thermostat (14.1), the bath warming to meet the
drink (14.2), and the room warming both (14.3). A full run strings them together as the bath moves from
regime to regime.

\end{document}
